\documentclass{article}
\usepackage{lmodern}
\usepackage{amsmath}
\usepackage{listings}
\usepackage{fullpage}
\usepackage{parskip}
\usepackage{xcolor}
\lstset{
	basicstyle=\ttfamily,
	columns=fullflexible,
	frame=single,
	breaklines=true,
	postbreak=\mbox{\textcolor{red}{$\hookrightarrow$}\space},
}
\begin{document}
\title{Experiment `p\_6\_4\_sum\_without\_smallest\_seq' Results}
\author{\today}
\date{}
\maketitle
\textbf{Experiment outcome: }SUCCESS
\\\textbf{Bad responses: }0
\\\textbf{Responses containing}~\texttt{assume}~\textbf{: }0
\\\textbf{Responses containing}~\texttt{expect}~\textbf{: }0
\\\textbf{Resolution attempts: }1
\\\textbf{Hard fails (resolution): }0
\\\textbf{Soft fails (resolution): }0
\\\textbf{Verification attempts: }1
\section*{Problem Specification}
\textbf{Problem name: }p\_6\_4\_sum\_without\_smallest\_seq
\\\textbf{Natural language statement: }Write a method that computes the sum of an array of values, except for the smallest one, in a single loop. In the loop, update the sum and the smallest value. After the loop, return the difference.
\\\textbf{Method signature: }p\_6\_4\_sum\_without\_smallest\_seq(inputs: seq<int>) returns (sum: int)
\subsection*{Ensures}
\begin{itemize}
\item \texttt{sum == sum\_array(inputs) - find\_smallest(inputs)}
\end{itemize}
\subsection*{Requires}
\begin{itemize}
\item \texttt{|inputs| >= 1}
\end{itemize}
\subsection*{Functional Code Given}
\begin{lstlisting}
function find_smallest(inputs: seq<int>): (ret: int)
  requires |inputs| >= 1
  ensures forall x :: x in inputs ==> ret <= x 
{
  if |inputs| == 1 then inputs[0]
  else 
    assert 2 <= |inputs|;
    var rest_smallest := find_smallest(inputs[1..]);
    var ret: int := if inputs[0] <= rest_smallest then inputs[0] else rest_smallest;
    assert inputs == [inputs[0]] + inputs[1..]; 
    assert ret == if inputs[0] <= rest_smallest then inputs[0] else rest_smallest;
    ret 
}

function sum_array(inputs: seq<int>): int
{
  if |inputs| == 0 then 0
  else  inputs[0] + sum_array(inputs[1..])
}
\end{lstlisting}
\clearpage
\section*{GenAI interactions}
Below you will find all interactions between the `user' (program) and the `assistant' (OpenAI).
\subsection*{Program $\rightarrow$ GenAI}
\begin{lstlisting}
You are given the following task to perform in Dafny:

Write a method that computes the sum of an array of values, except for the smallest one, in a single loop. In the loop, update the sum and the smallest value. After the loop, return the difference.

The signature should be:

p_6_4_sum_without_smallest_seq(inputs: seq<int>) returns (sum: int)

The method should respect the following contract:

requires |inputs| >= 1, ensures sum == sum_array(inputs) - find_smallest(inputs)

The contract uses the following dafny code:

function find_smallest(inputs: seq<int>): (ret: int)
  requires |inputs| >= 1
  ensures forall x :: x in inputs ==> ret <= x 
{
  if |inputs| == 1 then inputs[0]
  else 
    assert 2 <= |inputs|;
    var rest_smallest := find_smallest(inputs[1..]);
    var ret: int := if inputs[0] <= rest_smallest then inputs[0] else rest_smallest;
    assert inputs == [inputs[0]] + inputs[1..]; 
    assert ret == if inputs[0] <= rest_smallest then inputs[0] else rest_smallest;
    ret 
}

function sum_array(inputs: seq<int>): int
{
  if |inputs| == 0 then 0
  else  inputs[0] + sum_array(inputs[1..])
}

These function(s) must not be used in your implementation of the method.

Produce and show only the Dafny body of this method, including the curly braces that surround it. Do not show the signature nor contract. You must not use 'assume' nor 'expect' anywhere in your code.
\end{lstlisting}
\subsection*{GenAI $\rightarrow$ Program}
\textbf{System fingerprint: }null
\\\textbf{ID: }chatcmpl-Dn1VHBqfhXOwY3jUcSdq5XDfBGX1c
\\\textbf{Model: }gpt-5.5-2026-04-23
\\\textbf{Created at: }1780576223
\\\textbf{Finish reason: }stop
\\\textbf{Usage: }760 tokens in, and 6749 tokens out
\begin{lstlisting}
{
  var i: int := |inputs| - 1;
  assert 0 <= i < |inputs|;

  var runningSum: int := inputs[i];
  var smallest: int := inputs[i];

  assert |inputs[i..]| == 1;
  assert (inputs[i..])[0] == inputs[i];
  assert sum_array(inputs[i..]) == inputs[i];
  assert find_smallest(inputs[i..]) == inputs[i];

  while 0 < i
    invariant 0 <= i < |inputs|
    invariant runningSum == sum_array(inputs[i..])
    invariant smallest == find_smallest(inputs[i..])
    decreases i
  {
    var oldI: int := i;
    i := i - 1;

    assert oldI == i + 1;
    assert 0 <= i;
    assert i + 1 < |inputs|;
    assert runningSum == sum_array(inputs[i + 1..]);
    assert smallest == find_smallest(inputs[i + 1..]);

    assert |inputs[i..]| >= 2;
    assert (inputs[i..])[0] == inputs[i];
    assert (inputs[i..])[1..] == inputs[i + 1..];

    assert sum_array(inputs[i..]) == (inputs[i..])[0] + sum_array((inputs[i..])[1..]);
    assert sum_array(inputs[i..]) == inputs[i] + sum_array(inputs[i + 1..]);

    assert find_smallest(inputs[i..]) ==
      (if (inputs[i..])[0] <= find_smallest((inputs[i..])[1..]) then (inputs[i..])[0] else find_smallest((inputs[i..])[1..]));
    assert find_smallest(inputs[i..]) ==
      (if inputs[i] <= find_smallest(inputs[i + 1..]) then inputs[i] else find_smallest(inputs[i + 1..]));

    runningSum := inputs[i] + runningSum;
    assert runningSum == inputs[i] + sum_array(inputs[i + 1..]);
    assert runningSum == sum_array(inputs[i..]);

    if inputs[i] <= smallest {
      assert inputs[i] <= find_smallest(inputs[i + 1..]);
      smallest := inputs[i];
    } else {
      assert !(inputs[i] <= find_smallest(inputs[i + 1..]));
    }
    assert smallest ==
      (if inputs[i] <= find_smallest(inputs[i + 1..]) then inputs[i] else find_smallest(inputs[i + 1..]));
    assert smallest == find_smallest(inputs[i..]);
  }

  assert i == 0;
  assert inputs[i..] == inputs;
  assert runningSum == sum_array(inputs);
  assert smallest == find_smallest(inputs);
  sum := runningSum - smallest;
}
\end{lstlisting}
\section*{Final Program}
\begin{lstlisting}
function find_smallest(inputs: seq<int>): (ret: int)
  requires |inputs| >= 1
  ensures forall x :: x in inputs ==> ret <= x 
{
  if |inputs| == 1 then inputs[0]
  else 
    assert 2 <= |inputs|;
    var rest_smallest := find_smallest(inputs[1..]);
    var ret: int := if inputs[0] <= rest_smallest then inputs[0] else rest_smallest;
    assert inputs == [inputs[0]] + inputs[1..]; 
    assert ret == if inputs[0] <= rest_smallest then inputs[0] else rest_smallest;
    ret 
}

function sum_array(inputs: seq<int>): int
{
  if |inputs| == 0 then 0
  else  inputs[0] + sum_array(inputs[1..])
}

method p_6_4_sum_without_smallest_seq(inputs: seq<int>) returns (sum: int)
	requires |inputs| >= 1
	ensures sum == sum_array(inputs) - find_smallest(inputs)
{
  var i: int := |inputs| - 1;
  assert 0 <= i < |inputs|;

  var runningSum: int := inputs[i];
  var smallest: int := inputs[i];

  assert |inputs[i..]| == 1;
  assert (inputs[i..])[0] == inputs[i];
  assert sum_array(inputs[i..]) == inputs[i];
  assert find_smallest(inputs[i..]) == inputs[i];

  while 0 < i
    invariant 0 <= i < |inputs|
    invariant runningSum == sum_array(inputs[i..])
    invariant smallest == find_smallest(inputs[i..])
    decreases i
  {
    var oldI: int := i;
    i := i - 1;

    assert oldI == i + 1;
    assert 0 <= i;
    assert i + 1 < |inputs|;
    assert runningSum == sum_array(inputs[i + 1..]);
    assert smallest == find_smallest(inputs[i + 1..]);

    assert |inputs[i..]| >= 2;
    assert (inputs[i..])[0] == inputs[i];
    assert (inputs[i..])[1..] == inputs[i + 1..];

    assert sum_array(inputs[i..]) == (inputs[i..])[0] + sum_array((inputs[i..])[1..]);
    assert sum_array(inputs[i..]) == inputs[i] + sum_array(inputs[i + 1..]);

    assert find_smallest(inputs[i..]) ==
      (if (inputs[i..])[0] <= find_smallest((inputs[i..])[1..]) then (inputs[i..])[0] else find_smallest((inputs[i..])[1..]));
    assert find_smallest(inputs[i..]) ==
      (if inputs[i] <= find_smallest(inputs[i + 1..]) then inputs[i] else find_smallest(inputs[i + 1..]));

    runningSum := inputs[i] + runningSum;
    assert runningSum == inputs[i] + sum_array(inputs[i + 1..]);
    assert runningSum == sum_array(inputs[i..]);

    if inputs[i] <= smallest {
      assert inputs[i] <= find_smallest(inputs[i + 1..]);
      smallest := inputs[i];
    } else {
      assert !(inputs[i] <= find_smallest(inputs[i + 1..]));
    }
    assert smallest ==
      (if inputs[i] <= find_smallest(inputs[i + 1..]) then inputs[i] else find_smallest(inputs[i + 1..]));
    assert smallest == find_smallest(inputs[i..]);
  }

  assert i == 0;
  assert inputs[i..] == inputs;
  assert runningSum == sum_array(inputs);
  assert smallest == find_smallest(inputs);
  sum := runningSum - smallest;
}
\end{lstlisting}
\section*{Total Token Usage}
\textbf{Input tokens: }760
\\\textbf{Output tokens: }6749
\\\textbf{Reasoning tokens: }5997
\\\textbf{Sum of `total tokens': }7509
\section*{Experiment Timings}
\textbf{Overall Experiment} started at 1780576222306, ended at 1780576379202, lasting 156896ms (156.90 seconds)
\\\textbf{Iteration \#1} started at 1780576222311, ended at 1780576379202, lasting 156891ms (156.89 seconds)
\end{document}
